\documentclass[letterpaper]{article} % DO NOT CHANGE THIS
\usepackage[submission]{aaai2027}  % DO NOT CHANGE THIS
\usepackage[hyphens]{url}  % DO NOT CHANGE THIS
\usepackage{graphicx} % DO NOT CHANGE THIS
\urlstyle{rm} % DO NOT CHANGE THIS
\def\UrlFont{\rm}  % DO NOT CHANGE THIS
\usepackage{natbib}  % DO NOT CHANGE THIS AND DO NOT ADD ANY OPTIONS TO IT
\usepackage{caption} % DO NOT CHANGE THIS AND DO NOT ADD ANY OPTIONS TO IT
\frenchspacing  % DO NOT CHANGE THIS
\usepackage{float}

% Supplementary figures are numbered S1, S2, ... so they do not collide with the main paper.
\renewcommand{\thefigure}{S\arabic{figure}}

\pdfinfo{
/TemplateVersion (2027.1)
}

\title{Supplementary Material:\\ Attention Concentration Supports Transformer Length Generalization}
\author{}
\nocopyright

\begin{document}
\maketitle

This supplement collects the supporting figures for the co-varying variance statistic and the
variance-stabilizing intervention. Every numerical claim these figures illustrate is stated in the main
paper; the figures are provided here for completeness. Section and table references point to the main paper.

\begin{figure}[H]
\centering
\includegraphics[width=\columnwidth]{figures/fig_variance_collapse.pdf}
\caption{A co-varying statistic (main paper, Section~4.2).
Baseline attention-output variance in the deepest-collapsing layer, normalized to the training length,
falls with length on every task.}
\label{fig:var}
\end{figure}

\begin{figure}[H]
\centering
\includegraphics[width=\columnwidth]{figures/fig_prepost_var.pdf}
\caption{The intervention holds the statistic constant (main paper, Section~4.3).
For the intervention runs, the deep-layer variance before the normalization (solid) still collapses, while
the variance after it (dashed) is held near constant across length.}
\label{fig:prepost}
\end{figure}

\begin{figure}[H]
\centering
\includegraphics[width=\columnwidth]{figures/fig_accuracy_vs_length.pdf}
\caption{Per-token accuracy against length on a log axis, out to $50\times$ the training length (main paper,
Section~4.3).
Baseline is solid and the variance-stabilizing intervention is dashed, averaged over four seeds for the
retrieval tasks with min-to-max bands.
The intervention lies at or below the baseline in every retrieval panel.
Addition (bottom) fails from a short length in both conditions.}
\label{fig:acc}
\end{figure}

\begin{figure}[H]
\centering
\includegraphics[width=\columnwidth]{figures/fig_benefit_vs_length.pdf}
\caption{The intervention's per-token effect (intervention minus baseline) against length (main paper,
Section~4.3).
It is at most zero and becomes more negative with length, with the largest cost under RoPE.}
\label{fig:benefit}
\end{figure}

\begin{figure}[H]
\centering
\includegraphics[width=\columnwidth]{figures/fig_attn_base_vs_fix.pdf}
\caption{The behavior follows the operative variable (main paper, Section~4.4).
Attention on the correct source with the intervention off (solid) and on (dashed).
The intervention does not restore the attention, and under RoPE it lowers it, which matches the accuracy in
Table~1 of the main paper.}
\label{fig:attnfix}
\end{figure}

\end{document}
